\par\Needspace{6\baselineskip}\medskip
\begin{center}\hypertarget{native-173-article126}{}\pdfbookmark[2]{126.}{article-126}126.\end{center}
Das vorhergehende Verfahren führt unmittelbar zu einem \emph{dritten}. Wie zuvor bestimme man aus den Entfernungen des Himmelskörpers von der Erde bei der ersten und letzten Beobachtung die zugehörigen Längen in der Bahn samt den Radiusvektoren. Mithilfe der Lage der Bahnebene, die diese Rechnung liefert, ermittle man aus der mittleren Beobachtung die Länge in der Bahn und den Radiusvektor. Dann berechne man aus diesen drei heliozentrischen Orten die übrigen Elemente mit der Lösung der in den Artikeln 82 und 83 behandelten Aufgabe. Diese Rechnung ist von den Beobachtungszeiten unabhängig. So erhält man die drei mittleren Anomalien und die tägliche Bewegung, aus denen sich die Zeitabstände zwischen der ersten und zweiten Beobachtung sowie zwischen der zweiten und dritten berechnen lassen. Ihre Differenzen gegenüber den wahren Zeitabständen werden als $X$ und $Y$ genommen.

Dieses Verfahren wäre weniger geeignet, wenn die heliozentrische Bewegung nur einen kleinen Bogen umfasst. In einem solchen Fall hängt diese Bahnbestimmung, wie wir bereits in Artikel 82 bemerkt haben, von Größen dritter Ordnung ab und lässt daher keine hinreichende Genauigkeit zu. Geringste Änderungen der Werte von $x,y$ könnten sehr große Änderungen der Elemente und somit auch der Werte von $X,Y$ hervorrufen. Diese Änderungen dürfte man nicht als proportional zu den Änderungen der angesetzten Werte ansehen. Ist die heliozentrische Bewegung, die die drei Orte umfassen, jedoch beträchtlich, so wird das Verfahren gewiss sehr gut gelingen, sofern es durch die in den vorhergehenden Artikeln \SourceLeaf{174}erläuterten Ausnahmefälle nicht beeinträchtigt wird, auf die man offenbar auch hier achten muss.
